%% Chapter 13: Approximation Algorithms I
%% Status: drafted; figures are still placeholders.
%%
%% Written against slides 1-29 of "material/slides/Approximation Algorithms 2026.pdf"
%% (Johan van Rooij).  The pool exercise on asymmetric TSP is worked in the last
%% section before the exercises.
%%
%% Material deliberately NOT here because it is not in the deck: local search,
%% scheduling and load balancing, the primal-dual method.

\chapter{Approximation Algorithms I}
\label{ch:approx-i}

An approximation algorithm runs in polynomial time and returns a solution
together with a proof that the solution is not far from optimal. The proof is
the hard part, and the reason is worth saying at the outset: the guarantee
compares the algorithm's answer with an optimum that we cannot compute --- if
we could, we would not be approximating. Every argument in this chapter
therefore goes through some third quantity that we \emph{can} reason about, and
that is squeezed between the two.

\begin{goals}
  \poolgoal{define the approximation ratio for minimisation and for
        maximisation problems, and define the class APX}
  \poolgoal{give a $2$-approximation algorithm for \lang{Max SAT} and
        prove its ratio}
  \poolgoal{give the $2$-approximation and Christofides' algorithm for
        metric \lang{TSP}, and prove their ratios}
  \poolgoal{describe the general shape of an approximation ratio proof: a
        quantity that bounds both $\OPT$ and the value the algorithm returns}
  \poolgoal{give a $2$-approximation algorithm for \lang{Vertex Cover} and
        prove its ratio}
  \poolgoal{formulate the LP relaxation of \lang{Minimum Weight Vertex
        Cover}, give the rounding algorithm, and prove that it is a
        $2$-approximation}
  \poolgoal{explain the gap technique, and use it to show that \lang{TSP}
        without the triangle inequality admits no $c$-approximation for any
        constant $c$ unless $\mathrm{P} = \mathrm{NP}$}
  \poolgoal{give a polynomial-time $c$-approximation algorithm for
        \lang{TSP} instances that are not symmetric but satisfy the triangle
        inequality and $d(i,j) \le 2 d(j,i)$}
\end{goals}

%==============================================================================
\section{What an approximation guarantee is}
\label{sec:approx-definitions}
%==============================================================================

Everything we have done so far in this book has been aimed at one of two
targets: an optimal solution in polynomial time, or an optimal solution in
exponential time. For an \textrm{NP}-hard problem we can have either, but not
both. Once that is accepted, something has to give, and there are exactly three
things that can give: the running time (\cref{ch:exact-i,ch:exact-ii}), the
generality of the input (\cref{ch:fpt-i,ch:treewidth}), or the quality of the
answer. This chapter gives up on quality --- but not silently.

\begin{table}[ht]
\centering
\begin{tabular}{@{}lll@{}}
\toprule
 & \textbf{Optimal} & \textbf{Guaranteed quality} \\
\midrule
Polynomial time     & polynomial algorithms & \textbf{approximation algorithms} \\
Superpolynomial time & exact algorithms, ILP & hybrid algorithms \\
\bottomrule
\end{tabular}
\caption{Where approximation algorithms sit. The missing fourth entry ---
polynomial time, no guarantee --- is the world of construction heuristics,
local search and genetic algorithms, which are outside this course precisely
because they come with no proof.}
\label{tab:approx-landscape}
\end{table}

An approximation algorithm is a heuristic with a performance guarantee. Three
things are required of it.

\begin{enumerate}[leftmargin=*]
  \item It applies to an \emph{optimisation} problem, one where solutions have
        values and we want the largest or the smallest. Independent set, vertex
        cover, knapsack and \lang{Max SAT} qualify; \lang{SAT} and
        \lang{3-SAT} as decision problems do not, because there is nothing to
        be close to.
  \item It runs in polynomial time.
  \item It comes with a bound on the quality of what it returns.
\end{enumerate}

\begin{definition}[Approximation ratio]\label{def:approx-ratio}
Let $I$ be an instance of an optimisation problem, let $\OPT(I)$ be the value
of an optimal solution to $I$ and $\mathrm{ALG}(I)$ the value of the solution
returned by an algorithm. The algorithm has \emph{approximation ratio} $c$ if
for every instance $I$
\[
  \frac{\mathrm{ALG}(I)}{\OPT(I)} \le c \quad\text{(minimisation)},
  \qquad
  \frac{\OPT(I)}{\mathrm{ALG}(I)} \le c \quad\text{(maximisation)}.
\]
Such an algorithm is called a \emph{$c$-approximation algorithm}.
\end{definition}

The two cases are written so that the ratio is always at least $1$, with $1$
meaning exactly optimal and larger meaning worse. That convention costs one
moment of care --- for a maximisation problem the optimum is on top --- and
saves a great deal of confusion later, because it makes ``ratio $2$'' mean the
same thing everywhere.

\begin{definition}[APX]\label{def:apx}
\textrm{APX} is the class of optimisation problems in \textrm{NP} that admit a
polynomial-time $c$-approximation algorithm for some constant $c$.
\end{definition}

A $c$-approximation with $c$ constant is also called a \emph{constant factor}
approximation. \Cref{ch:approx-ii} will refine \cref{def:apx} into a scale of
five different qualities of approximation, and \cref{ch:hardness} will show how
to prove that a problem is not in \textrm{APX}. For this chapter, constant
factors are the whole story.

%==============================================================================
\section{Maximum satisfiability}
\label{sec:max-sat}
%==============================================================================

The first example is chosen because its algorithm is a single line and its
analysis is two.

\begin{probdef}{Max SAT}
  \instance A set $C$ of clauses over a set $X$ of Boolean variables, and an
  integer $k$.
  \question Is there a truth assignment to $X$ satisfying at least $k$ clauses
  of $C$?
\end{probdef}

As an optimisation problem, we are asked for an assignment satisfying as many
clauses as possible. \lang{Max $k$-SAT} is the same problem with every clause
restricted to at most $k$ literals. Note that \lang{Max SAT} is interesting
even on formulas that are not satisfiable, which is exactly what makes it an
optimisation problem and \lang{SAT} not one.

\begin{algorithm}[ht]
\caption{$2$-approximation for \lang{Max SAT}}
\label{alg:maxsat}
\begin{algorithmic}[1]
  \State let $a$ be any truth assignment to $X$
  \State let $\bar a$ be $a$ with every variable negated
  \State \Return whichever of $a$ and $\bar a$ satisfies more clauses
\end{algorithmic}
\end{algorithm}

\begin{theorem}\label{thm:maxsat}
\Cref{alg:maxsat} is a $2$-approximation for \lang{Max SAT}.
\end{theorem}

\begin{proof}
Every clause contains at least one literal, and a literal is satisfied by
exactly one of $a$ and $\bar a$. So every clause is satisfied by at least one
of the two assignments, and the two clause counts sum to at least $|C|$. The
better of the two therefore satisfies at least $|C|/2$ clauses, so
$\mathrm{ALG} \ge |C|/2$. On the other side no assignment can satisfy more
clauses than there are, so $\OPT \le |C|$. Hence
$\OPT/\mathrm{ALG} \le |C| / (|C|/2) = 2$.
\end{proof}

Notice the shape of that proof, because it is the shape of every proof in this
chapter: a quantity --- here the number of clauses $|C|$ --- that bounds $\OPT$
from one side and $\mathrm{ALG}$ from the other. We return to this in
\cref{sec:approx-proof-shape}.

%==============================================================================
\section{Metric TSP}
\label{sec:metric-tsp}
%==============================================================================

The second example is the problem approximation algorithms are most often
illustrated with, and the one where the choice of the third quantity is
prettiest.

\begin{probdef}{TSP}
  \instance A set of $n$ cities and a distance $d(x,y)$ for every pair.
  \question A shortest cycle visiting every city exactly once.
\end{probdef}

In general this problem is hopeless to approximate, as \cref{sec:gap-technique}
will show. Two restrictions rescue it, and both are natural for distances that
come from an actual map:
\begin{itemize}[nosep,leftmargin=*]
  \item \emph{symmetry}: $d(x,y) = d(y,x)$ for all $x,y$;
  \item \emph{the triangle inequality}: $d(x,y) + d(y,z) \ge d(x,z)$ for all
        $x,y,z$, that is, a detour never helps.
\end{itemize}
An instance satisfying both is called \emph{metric}. Assume both from here to
the end of the section.

\subsection{A spanning tree is a lower bound}

The whole of metric \lang{TSP} approximation rests on one observation.

\begin{lemma}\label{lem:tsp-mst-lower-bound}
Let $\mathrm{MST}$ be the weight of a minimum spanning tree of the instance.
Then $\OPT \ge \mathrm{MST}$.
\end{lemma}

\begin{proof}
Delete any one edge from an optimal tour. What remains is a Hamiltonian path,
which is in particular a spanning tree, and its weight is at most that of the
tour since weights are non-negative. A minimum spanning tree is no heavier.
\end{proof}

That is the third quantity. It bounds $\OPT$ from below, and the algorithm will
be built so that its output is bounded from above by twice the same quantity.

\begin{algorithm}[ht]
\caption{Double-tree $2$-approximation for metric \lang{TSP}}
\label{alg:tsp-double-tree}
\begin{algorithmic}[1]
  \State compute a minimum spanning tree $T$
  \State root $T$ anywhere and list its vertices in preorder, visiting a node
         before its children
  \State \Return the tour that visits the cities in that order
\end{algorithmic}
\end{algorithm}

\begin{theorem}\label{thm:tsp-double-tree}
\Cref{alg:tsp-double-tree} is a $2$-approximation for metric \lang{TSP}.
\end{theorem}

\begin{proof}
Walking around the tree, traversing every edge once downwards and once
upwards, is a closed walk of weight exactly $2\,\mathrm{MST}$ visiting every
city. The preorder tour is that walk with repeated visits shortcut away, and by
the triangle inequality every shortcut only shortens it. Hence
$\mathrm{ALG} \le 2\,\mathrm{MST}$. With \cref{lem:tsp-mst-lower-bound},
\[ \frac{\mathrm{ALG}}{\OPT} \le \frac{2\,\mathrm{MST}}{\mathrm{MST}} = 2. \]
\end{proof}

\begin{figure}[ht]
\centering
\figplaceholder{5cm}{A minimum spanning tree of a small point set; the closed
walk that traverses every tree edge twice; and the tour obtained by shortcutting
repeated visits. The shortcuts are the only place the triangle inequality is
used.}
\caption{The double-tree algorithm. Doubling the tree makes every degree even,
so a closed walk exists; shortcutting turns the walk into a tour and, by the
triangle inequality, cannot make it longer.}
\label{fig:tsp-double-tree}
\end{figure}

\subsection{Christofides' algorithm}

Doubling the tree is wasteful. The only reason to double was to make every
degree even, and most vertices of a tree already have even degree. Repairing
just the odd ones is cheaper, and this is Christofides' idea.

\begin{algorithm}[ht]
\caption{Christofides' $3/2$-approximation for metric \lang{TSP}}
\label{alg:christofides}
\begin{algorithmic}[1]
  \State compute a minimum spanning tree $T$
  \State let $W = \{v : v \text{ has odd degree in } T\}$
  \State compute a minimum weight perfect matching $M$ in the complete graph
         induced on $W$
  \State form the multigraph $T + M$, which is Eulerian
  \State compute an Euler tour $C'$ of $T + M$
  \State \Return the tour obtained from $C'$ by shortcutting repeated visits
\end{algorithmic}
\end{algorithm}

Three details deserve a word. First, $|W|$ is even, because in any graph the
number of odd-degree vertices is even, so the matching exists. Second, adding a
perfect matching on $W$ raises the degree of exactly the odd vertices by one, so
every vertex of $T+M$ has even degree and the multigraph is connected: it is
Eulerian --- and a connected multigraph all of whose degrees are even has a
closed walk using every edge exactly once, found in linear time by repeatedly
following unused edges and splicing in the detours. Third, the
matching is computed by the algorithm of \cref{ch:matching}, in polynomial
time --- this step is the reason Christofides' algorithm is not a five-line
exercise.

\begin{theorem}\label{thm:christofides}
\Cref{alg:christofides} is a $3/2$-approximation for metric \lang{TSP}.
\end{theorem}

\begin{proof}
The weight of $T$ is at most $\OPT$ by \cref{lem:tsp-mst-lower-bound}.

For the matching, consider an optimal tour and restrict it to the vertices of
$W$: shortcutting the tour onto $W$ gives a cycle through the vertices of $W$
of weight at most $\OPT$, again by the triangle inequality. That cycle has even
length, so its edges split into two perfect matchings on $W$, and the lighter of
the two has weight at most $\OPT/2$. A minimum weight perfect matching is no
heavier, so $w(M) \le \OPT/2$.

Hence $w(T+M) \le \OPT + \OPT/2 = \tfrac32 \OPT$. The Euler tour has exactly
that weight, and shortcutting does not increase it, so
$\mathrm{ALG} \le \tfrac32\OPT$.
\end{proof}

\begin{figure}[ht]
\centering
\figplaceholder{5.5cm}{Four panels: the minimum spanning tree $T$; the odd
degree vertices $W$ marked; the minimum weight perfect matching $M$ on $W$ drawn
in a second colour; and the Eulerian multigraph $T+M$ with its Euler tour.}
\caption{Christofides' algorithm. Only the odd-degree vertices need repairing,
and repairing them costs at most half an optimal tour.}
\label{fig:christofides}
\end{figure}

\begin{intuition}
The step to watch is the bound on the matching. It is the only place where the
optimal tour is used as a \emph{source} of a cheap object rather than as a
target to compare against: an optimal tour, restricted to $W$, is itself a
union of two matchings, so one of them is cheap, so the minimum one is cheap.
This trick --- get an upper bound on something we compute by cutting up an
object we cannot compute --- recurs throughout approximation.
\end{intuition}

%==============================================================================
\section{The shape of an approximation proof}
\label{sec:approx-proof-shape}
%==============================================================================

Both TSP proofs, and the \lang{Max SAT} proof before them, have the same three
lines. It is worth isolating the pattern, because most approximation ratio
proofs in this course and beyond follow it.

\begin{enumerate}[leftmargin=*]
  \item Find a quantity $x$ that relates to both the optimum and the output.
        For metric \lang{TSP} it was the weight of a minimum spanning tree, for
        \lang{Max SAT} the number of clauses.
  \item Prove $\OPT \ge c_1 x$ (for a minimisation problem; $\OPT \le c_1 x$
        for maximisation).
  \item Prove $\mathrm{ALG} \le c_2 x$.
  \item Combine: the ratio is at most $c_2/c_1$.
\end{enumerate}

The art is entirely in step~1. Steps~2 and~3 are usually short once $x$ is
right, and hopeless when it is wrong. A good $x$ is something the algorithm
actually computes --- that makes step~3 nearly free --- and something an
optimal solution visibly contains --- that makes step~2 nearly free. A minimum
spanning tree is both.

In \cref{sec:lp-rounding} we will see the systematic way of producing such an
$x$: solve a relaxation of the problem, one whose optimum is by construction a
bound on $\OPT$ and which we can compute in polynomial time.

%==============================================================================
\section{The gap technique}
\label{sec:gap-technique}
%==============================================================================

Approximation ratios are proved; the absence of approximation ratios has to be
proved too, and the basic tool is a hardness reduction that is deliberately
wasteful.

\begin{definition}[Gap technique]\label{def:gap-technique}
A \emph{gap-creating reduction} is an \textrm{NP}-hardness proof whose output
instances have optimum value at most $a$ when the input is a yes-instance and
at least $b > a$ when it is a no-instance. An approximation algorithm with
ratio better than $b/a$ would then distinguish the two cases, and hence solve
the \textrm{NP}-hard problem in polynomial time.
\end{definition}

The whole content is in making $b/a$ large. Since the two cases must be
separated by a wide margin, the reduction deliberately blows up the cost of the
no-instances.

\begin{theorem}\label{thm:tsp-no-approx}
Unless $\mathrm{P} = \mathrm{NP}$, no polynomial-time algorithm approximates
\lang{TSP} without the triangle inequality within any constant ratio $c > 1$.
\end{theorem}

\begin{proof}
Let $c > 1$ be given, and let $G = (V,E)$ be an instance of
\lang{Hamiltonian Cycle} with $n = |V|$. Build a \lang{TSP} instance with one
city per vertex and
\[
  d(v,w) = \begin{cases}
    1 & \text{if } \{v,w\} \in E, \\
    nc + 1 & \text{otherwise.}
  \end{cases}
\]
If $G$ has a Hamiltonian cycle, that cycle is a tour of length exactly $n$. If
$G$ has none, every tour must use at least one non-edge, so its length is at
least $(n-1) + nc + 1 > nc$.

The two cases are separated by a factor of more than $c$. A $c$-approximation
algorithm run on the constructed instance returns a tour of length at most
$c \cdot \OPT$, which in the first case is at most $cn$ and in the second is
more than $cn$; comparing the returned length with $cn$ therefore decides
\lang{Hamiltonian Cycle} in polynomial time.
\end{proof}

Two remarks. First, the constructed instance badly violates the triangle
inequality --- a non-edge of length $nc+1$ can be short-circuited by a path of
two edges of length $1$ --- which is exactly why \cref{thm:christofides} is not
contradicted. Second, the argument works for any $c$, even one growing with
$n$; \cref{ch:approx-ii} pushes it to polynomial $c$.

%==============================================================================
\section{Vertex cover}
\label{sec:vertex-cover-approx}
%==============================================================================

The next example is the one everybody meets first, and its analysis is three
lines. What is worth watching is where the lower bound on $\OPT$ comes from:
the algorithm produces it as a by-product of doing its job.

\begin{probdef}{Vertex Cover}
  \instance An undirected graph $G = (V,E)$.
  \question A smallest set $C \subseteq V$ such that every edge has at least
  one endpoint in $C$.
\end{probdef}

\begin{algorithm}[ht]
\caption{$2$-approximation for \lang{Vertex Cover}}
\label{alg:vertex-cover}
\begin{algorithmic}[1]
  \State $E' \gets E$, \quad $C \gets \emptyset$
  \While{$E' \ne \emptyset$}
    \State pick any edge $\{u,v\} \in E'$
    \State $C \gets C \cup \{u,v\}$
    \State remove from $E'$ every edge incident to $u$ or to $v$
  \EndWhile
  \State \Return $C$
\end{algorithmic}
\end{algorithm}

The algorithm clearly runs in polynomial time and clearly returns a vertex
cover: an edge leaves $E'$ only when one of its endpoints enters $C$.

\begin{theorem}\label{thm:vertex-cover-approx}
\Cref{alg:vertex-cover} is a $2$-approximation for \lang{Vertex Cover}.
\end{theorem}

\begin{proof}
Let $A$ be the set of edges picked in the loop. No two of them share an
endpoint, since picking an edge removes every edge incident to either endpoint;
so $A$ is a matching. Any vertex cover must contain at least one endpoint of
each edge of $A$, and these endpoints are distinct, so $\OPT \ge |A|$. The
algorithm returns exactly the $2|A|$ endpoints, so
$\mathrm{ALG} = 2|A| \le 2\OPT$.
\end{proof}

Here the third quantity $x$ is $|A|$, the size of the matching the algorithm
happens to build. Note that the algorithm is greedy in a very weak sense --- it
picks \emph{any} edge, not a cleverly chosen one --- and that the analysis does
not care. Being unable to improve on this by any clever choice of edge is not
an accident: \cref{ch:hardness} shows that \lang{Vertex Cover} is
\textrm{APX}-hard, so no polynomial-time algorithm has ratio $1 + \varepsilon$
for every $\varepsilon$ unless $\mathrm{P} = \mathrm{NP}$.

%==============================================================================
\section{Minimum weight vertex cover by LP rounding}
\label{sec:lp-rounding}
%==============================================================================

Now give each vertex $v$ a weight $w_v \ge 0$ and ask for a vertex cover of
minimum total weight. \Cref{alg:vertex-cover} is no longer a $2$-approximation,
and the reason is plain: picking an arbitrary edge may put two very heavy
vertices into $C$ when one light vertex would have covered everything.

What saves us is a general method. Write the problem as an integer linear
program, drop the integrality, solve what remains, and round.

\begin{definition}[ILP formulation]\label{def:vc-ilp}
With one variable $x_v$ per vertex, \lang{Minimum Weight Vertex Cover} is
\[
\begin{array}{lll}
  \text{minimise}   & \displaystyle\sum_{v \in V} w_v x_v & \\[2pt]
  \text{subject to} & x_u + x_v \ge 1 & \text{for each } \{u,v\} \in E, \\[2pt]
                    & x_v \in \{0,1\} & \text{for each } v \in V.
\end{array}
\]
The \emph{LP relaxation} replaces the last line by $0 \le x_v \le 1$.
\end{definition}

Integer linear programming is \textrm{NP}-hard --- of course it is, we have
just encoded an \textrm{NP}-hard problem in it. Linear programming is not: it
is solvable in polynomial time by the ellipsoid method or by interior point
methods. (Not by the simplex algorithm, which is fast in practice but takes
exponential time in the worst case.) So the relaxation is something we can
actually compute, and since every vertex cover is a feasible solution of the
relaxation, its optimum $z^{*}$ satisfies $z^{*} \le \OPT$. That is the lower
bound the previous sections had to find by hand.

\begin{algorithm}[ht]
\caption{LP rounding for \lang{Minimum Weight Vertex Cover}}
\label{alg:vc-lp-rounding}
\begin{algorithmic}[1]
  \State solve the LP relaxation of \cref{def:vc-ilp}, giving $x^{*}$
  \State \Return $C = \{v : x^{*}_v \ge 1/2\}$
\end{algorithmic}
\end{algorithm}

\begin{theorem}\label{thm:vc-lp-rounding}
\Cref{alg:vc-lp-rounding} runs in polynomial time and is a $2$-approximation
for \lang{Minimum Weight Vertex Cover}.
\end{theorem}

\begin{proof}
It returns a vertex cover: for every edge $\{u,v\}$ the constraint
$x^{*}_u + x^{*}_v \ge 1$ forces at least one of the two to be at least $1/2$.

For the ratio, let $z^{*} = \sum_v w_v x^{*}_v$ be the LP optimum. Every vertex
cover gives a feasible integral solution of the LP, so $z^{*} \le \OPT$. In the
other direction, discarding the vertices with $x^{*}_v < 1/2$ only decreases the
sum, and on the vertices that remain $x^{*}_v \ge 1/2$, so
\[
  z^{*} \;=\; \sum_{v \in V} w_v x^{*}_v
        \;\ge\; \sum_{v \,:\, x^{*}_v \ge 1/2} w_v x^{*}_v
        \;\ge\; \frac12 \sum_{v \,:\, x^{*}_v \ge 1/2} w_v
        \;=\; \frac12\,\mathrm{ALG}.
\]
Hence $\mathrm{ALG} \le 2 z^{*} \le 2\OPT$.
\end{proof}

\begin{intuition}
The rounding threshold $1/2$ is not a tuning parameter; it is forced. The
constraint has two variables summing to at least $1$, so $1/2$ is the largest
threshold that is guaranteed to catch one of them, and it is also the loss
factor in the second half of the proof. A constraint with $k$ variables summing
to at least $1$ would give threshold $1/k$ and ratio $k$ by the identical
argument --- which is how the same method approximates set cover when every
element lies in at most $k$ sets.
\end{intuition}

The general method behind \cref{alg:vc-lp-rounding} is worth stating on its
own, since it is one of the two or three most useful things in the subject.

\begin{enumerate}[leftmargin=*]
  \item Write the problem as an integer program.
  \item Relax the integrality constraints and solve the resulting LP in
        polynomial time. Its value $z^{*}$ is a bound on $\OPT$, for free.
  \item Round the fractional solution to an integral one, and bound how much
        the rounding costs.
\end{enumerate}
Step~3 is where the thinking happens, and the ratio one gets is exactly the
worst-case cost of the rounding --- what is called the \emph{integrality gap}
of the formulation.

%==============================================================================
\section{Asymmetric instances}
\label{sec:asymmetric-tsp}
%==============================================================================

Real distances are often not symmetric. An electric car in the mountains uses
more battery going up than coming down, so $d(i,j) \ne d(j,i)$, and neither
\cref{alg:tsp-double-tree} nor \cref{alg:christofides} applies as it stands ---
a spanning tree has no direction, and shortcutting a walk uses the triangle
inequality in both directions.

If the asymmetry is bounded, however, the whole analysis survives at a constant
loss. This is the pool exercise \emph{Approximation non-symmetric TSP}, and it
is worth working through because the
argument is a template for what to do with a hypothesis of the form ``almost
symmetric''.

\begin{theorem}\label{thm:asymmetric-tsp}
Consider \lang{TSP} instances that satisfy the triangle inequality
$d(i,j) + d(j,k) \ge d(i,k)$ and the bounded asymmetry condition
$d(i,j) \le 2\,d(j,i)$ for all cities $i,j,k$. There is a polynomial-time
$3$-approximation algorithm for these instances.
\end{theorem}

\begin{proof}
Symmetrise the instance by taking, for every pair, the \emph{larger} of the two
directions:
\[ \hat d(i,j) \;=\; \hat d(j,i) \;=\; \max\{d(i,j),\, d(j,i)\}. \]
Run Christofides' algorithm on $\hat d$ and orient the tour it returns
arbitrarily. That is the whole algorithm; it is polynomial because
\cref{alg:christofides} is. Three claims prove the ratio.

\emph{Claim 1: $\hat d$ is a metric.} It is symmetric and non-negative by
construction. For the triangle inequality, fix $i,j,k$. By definition
$\hat d(i,k)$ is either $d(i,k)$ or $d(k,i)$. In the first case
\[ d(i,k) \le d(i,j) + d(j,k) \le \hat d(i,j) + \hat d(j,k), \]
and in the second case, using the triangle inequality along $k \to j \to i$,
\[ d(k,i) \le d(k,j) + d(j,i) \le \hat d(j,k) + \hat d(i,j). \]
Either way $\hat d(i,k) \le \hat d(i,j) + \hat d(j,k)$, so
\cref{alg:christofides} applies to $\hat d$ and returns a tour of $\hat
d$-length at most $\tfrac32 \widehat{\OPT}$, where $\widehat{\OPT}$ is the
optimum of the symmetrised instance.

\emph{Claim 2: $\widehat{\OPT} \le 2\OPT$.} This is the only place the bounded
asymmetry is used. Take an optimal directed tour and forget the directions. For
each of its arcs $(i,j)$ we have $d(j,i) \le 2 d(i,j)$ by hypothesis, hence
$\hat d(i,j) = \max\{d(i,j),d(j,i)\} \le 2 d(i,j)$. Summing over the arcs, the
undirected tour has $\hat d$-length at most $2\OPT$, and $\widehat{\OPT}$ is no
larger.

\emph{Claim 3: orienting costs nothing.} Let $T$ be the undirected tour
returned, oriented arbitrarily. Every arc of the oriented tour satisfies
$d(i,j) \le \max\{d(i,j),d(j,i)\} = \hat d(i,j)$, so the directed length of $T$
is at most its $\hat d$-length.

Chaining the three claims,
\[ d(T) \;\le\; \hat d(T) \;\le\; \tfrac32 \widehat{\OPT}
        \;\le\; \tfrac32 \cdot 2\OPT \;=\; 3\OPT. \qedhere \]
\end{proof}

Using \cref{alg:tsp-double-tree} instead of Christofides' algorithm gives the
same argument with $\tfrac32$ replaced by $2$, hence a $4$-approximation, and
avoids computing a matching. Both are acceptable answers to the pool question,
which asks only for some constant.

\begin{intuition}
Why the \emph{maximum} and not the minimum? Taking the minimum makes the
symmetrised distances smaller, which looks better --- but it destroys the
triangle inequality, and then Christofides' analysis no longer applies. Taking
the maximum keeps the triangle inequality for free (claim~1 does not use the
bounded asymmetry at all) and pays for it in claim~2, where the hypothesis
$d(i,j) \le 2 d(j,i)$ says exactly that the maximum is never more than twice
the direction an optimal tour actually travels. The hypothesis is used once,
and it buys one factor of two.
\end{intuition}

The point of the exercise is not the constant. It is that a bounded asymmetry
can be absorbed into the constant at all: the instance is turned into one the
existing algorithms accept, and each step of the translation is charged
separately.

%==============================================================================
\section{Notes and further reading}
\label{sec:approx-i-further-reading}
%==============================================================================

This chapter follows Johan van Rooij's lecture slides. For the design and
analysis of approximation algorithms, \citet[Chapter~35]{CLRS2022} covers
vertex cover, TSP and set cover at the level of this chapter;
\citet{WilliamsonShmoys2011} is the standard modern text, organised by
technique rather than by problem, and \citet{Vazirani2001} is the other
standard. For the landscape of approximability, which is the subject of
\cref{ch:approx-ii}, \citet[Chapter~3]{Ausiello1999} is the reference.

Christofides' algorithm was a 1976 Carnegie Mellon technical report that was
never submitted anywhere and accumulated several thousand citations in that
state; it finally appeared in a journal in \citet{Christofides2022}. Its ratio
stood as the best known for metric TSP for forty-five years, until
\citet{KarlinKleinOveisGharan2021} obtained $3/2 - \varepsilon$ for an
$\varepsilon$ of about $10^{-36}$. The LP rounding argument for weighted vertex cover is due
to \citet{Hochbaum1982} and \citet{BarYehudaEven1981}.

%==============================================================================
\section{Exercises}
\label{sec:ex-approx-i}
%==============================================================================

\begin{exercise}
Give a $2$-approximation for \lang{Max $k$-SAT} for every $k$, and improve it
to a $(1 - 2^{-k})^{-1}$-approximation. Hint: for the second part, count how
many of the $2^k$ assignments of the variables in one clause fail to satisfy it.
\end{exercise}

\begin{exercise}
Show that \cref{lem:tsp-mst-lower-bound} fails without the assumption that
distances are non-negative, and that the double-tree algorithm then has no
constant ratio.
\end{exercise}

\begin{exercise}
Find an instance of metric \lang{TSP} on which \cref{alg:tsp-double-tree}
returns a tour of length close to $2\OPT$, showing that the analysis of
\cref{thm:tsp-double-tree} is tight.
\end{exercise}

\begin{exercise}
In Christofides' algorithm, why is it not enough to compute \emph{any} perfect
matching on $W$? Where exactly does the proof of \cref{thm:christofides} use
minimality?
\end{exercise}

\begin{exercise}
The greedy algorithm that repeatedly picks a vertex of maximum degree is a
plausible alternative to \cref{alg:vertex-cover}. Show that it does not have
approximation ratio $2$.
\end{exercise}

\begin{exercise}
Show that the LP relaxation of \cref{def:vc-ilp} can have optimum $|V|/2$ on a
graph whose minimum vertex cover has size $|V| - 1$. Hint: a complete graph.
Conclude that no rounding whatsoever can turn this formulation into a
$(2-\varepsilon)$-approximation for a fixed $\varepsilon > 0$.
\end{exercise}

\begin{exercise}
Apply the gap technique to \lang{Chromatic Number}: use the
\textrm{NP}-hardness of deciding $3$-colourability to show that no
polynomial-time algorithm approximates the chromatic number within a ratio
better than $4/3$ unless $\mathrm{P} = \mathrm{NP}$.
\end{exercise}

\begin{exercise}
In \cref{thm:asymmetric-tsp}, where is the hypothesis $d(i,j) \le 2 d(j,i)$
used, and what ratio does the same proof give if the hypothesis is weakened to
$d(i,j) \le \alpha\, d(j,i)$ for a constant $\alpha$?
\end{exercise}
